
\subsubsection{Controle das juntas do manipulador robótico}
\label{sec:controle_PD_juntas}

O \gls{mr} acoplado à base $\{B\}$ possui 3 graus de liberdade rotacionais, descritos pelas coordenadas generalizadas

\begin{equation}
\vec\theta =
\begin{bmatrix}
\theta_1 \\ \theta_2 \\ \theta_3
\end{bmatrix},
\end{equation}

e por suas respectivas velocidades articulares $\dot{\vec\theta}$.
As referências articulares são denotadas por

\begin{equation}
\vec\theta_{\text{ref}} =
\begin{bmatrix}
\theta_{1,\text{ref}} \\
\theta_{2,\text{ref}} \\
\theta_{3,\text{ref}}
\end{bmatrix},
\end{equation}

que são obtidas a partir da solução do problema de cinemática inversa, conforme a Seção~\eqref{sec:cinematica_inversa}, do efetuador para uma trajetória desejada de aproximação e alinhamento com a porta de acoplamento.
Define-se o erro articular e sua derivada como

\begin{equation}
\vec e_\theta = \vec\theta - \vec\theta_{\text{ref}},
\qquad
\dot{\vec e}_\theta = \dot{\vec\theta} - \dot{\vec\theta}_{\text{ref}},
\end{equation}

em que as referências de velocidade $\dot{\vec\theta}_{\text{ref}}$ são nulas ou variam lentamente, caracterizando trajetórias suavizadas para o manipulador.
A lei de controle PD em espaço de juntas é então escrita como

\begin{equation}
\vec\tau_\theta
=
- K_{p,\theta}\,\vec e_\theta
- K_{d,\theta}\,\dot{\vec e}_\theta,
\label{eq:controle_PD_juntas}
\end{equation}

onde $\vec\tau_\theta$
é o vetor de torques nas juntas,
e $K_{p,\theta}$, $K_{d,\theta}$
são matrizes com os ganhos de posição e velocidade, respectivamente.
Neste contexto, o \gls{mr} é controlado no espaço de juntas \cite{craig2005}, enquanto o vínculo com o espaço operacional (posição e orientação do efetuador em relação à porta de acoplamento) é realizado pela cinemática direta e inversa discutidas anteriormente.
